% !TeX root = ../../Book3.tex
% 纲要标题：### 20.3 高度一局部环与除子类群
% 纲要位置：计划纲要.md，第 2598 行。

\section{高度一局部环与除子类群}
\label{b3:sec:2003}

正规整环的高度一局部环是 DVR，每个这样的素理想因此给出一个离散阶数。这些阶数可以检测分式是否属于原环，也能记录主理想不能表示的除子信息。


% 纲要内容（仅作写作计划，尚非教材正文）：
%
% * 正规整环是分式域中全部高度一局部环的交；用负赋值检测分式不属于环
% * 归一化离散赋值；正阶对应主理想的极小素理想，故主除子支撑有限
% * Weil 除子作为有限整数阶数资料，类群为零回用高度一素理想主性判据
% * 函数与普通理想的高度一恢复不同；类群局部化为删除素理想类所得的商群
%

\subsection{正规整环与高度一局部化}
\label{b3:subsec:2003:01}

\begin{theorem}\label{b3:thm:normal-height-one-intersection}
设 \(R\) 为交换 Noether 正规整环，\(K=\operatorname{Frac}(R)\)。则在 \(K\) 中有
\[
R=\bigcap_{\operatorname{ht}\mathfrak p=1}R_{\mathfrak p},
\]
其中交集遍历 \(R\) 的所有高度一素理想 \(\mathfrak p\)，每个 \(R_{\mathfrak p}\) 都是 DVR。若没有高度一素理想，则 \(R\) 为域，交集按 \(K\) 解释。
\end{theorem}
\begin{proof}
DVR 结论由正规性及\cref{b3:thm:dvr-characterizations}得到。包含关系从左到右显然；反向把交集中的元素写成 \(b/a\)，其中 \(0\ne a\in R\)。Serre 判据使 \(R\) 满足 \(S_2\)，所以\cref{b3:lem:s2-fraction-membership}使 \(b\in aR\)。若无高度一素理想，任意非零非单位的主理想都应有一个高度一极小素理想，故不存在这样的非单位，\(R\) 是域。
\end{proof}

记 \(v_{\mathfrak p}:K^*\rightarrow\mathbb Z\) 为 \(R_{\mathfrak p}\) 的归一化离散赋值，即统一化元的值为 \(1\)。对 \(0\ne f\in K\)，上述定理等价于
\[
f\in R\quad\Longleftrightarrow\quad v_{\mathfrak p}(f)\ge0
\quad\text{对所有高度一}\;\mathfrak p.
\]
并且 \(f\in R^*\) 当且仅当所有赋值均为零，因为这时 \(f\) 与 \(f^{-1}\) 同时属于 \(R\)。环元素与单位因此可由各个余维一处的阶数判定。从函数的角度看，一个分式若在所有余维一处都没有极点，就已经属于原环，无须再到余维至少二的点作额外检测。正规性使函数的高余维延拓由这些阶数控制。

\ProseParagraphBreak
例如，设 \(k\) 为域，\(R=k[x,y]\)。在 \(R_{(y)}\) 中，\(y\) 是统一化元而 \(x\) 是单位，所以
\[
v_{(y)}(x/y)=0-1=-1.
\]
仅这一个高度一素理想便检测到 \(x/y\notin R\)。\((x,y)\) 的高度为二，虽然它在几何上是一个特殊的闭点，却不是这里给出离散阶数的位置。高余维处的理想信息并未消失；对有理函数的环元素判定，高度一检查已经足够。

\subsection{Weil 除子与类群}
\label{b3:subsec:2003:02}

每个高度一素理想给出一个整数阶数。把仅有有限多处非零的整数分别指定给这些素理想，就得到一个 Weil 除子；它可以看成有限支撑的阶数资料。同一个非零分式在各处产生的阶数给出主除子。类群所比较的，正是任意这样的整数资料与一个全局分式能够同时产生的资料。

\begin{definition}\label{b3:def:weil-divisor-class-group}
设 \(R\) 为交换 Noether 正规整环，\(K=\operatorname{Frac}(R)\)。对每个高度一素理想 \(\mathfrak p\)，令 \(v_{\mathfrak p}:K^*\to\mathbb Z\) 为离散赋值环 \(R_{\mathfrak p}\) 的归一化离散赋值，即统一化元的值为 \(1\)。以高度一素理想为基的自由交换群
\[
\operatorname{Div}(R)=\bigoplus_{\operatorname{ht}\mathfrak p=1}\mathbb Z[\mathfrak p]
\]
称为 \(R\) 的\term[Weilchuziqun]{Weil 除子群}[Weil divisor group]。其元素 \(D=\sum n_{\mathfrak p}[\mathfrak p]\) 为有限和，称为 Weil 除子；若各 \(n_{\mathfrak p}\ge0\)，则称 \(D\) 为有效除子。对 \(f\in K^*\)，定义主除子
\[
\operatorname{div}(f)=\sum_{\operatorname{ht}\mathfrak p=1}v_{\mathfrak p}(f)[\mathfrak p].
\]
商群
\(\operatorname{Cl}(R)=\operatorname{Div}(R)/\{\operatorname{div}(f):f\in K^*\}\)
称为\term[chuzileiqun]{除子类群}[divisor class group]。两个除子之差为主除子时，称它们线性等价。
\end{definition}
\symidx{\(\operatorname{Div}(R)\)：正规整环的 Weil 除子群}
\symidx{\(\operatorname{div}(f)\)：非零分式 \(f\) 的主除子}
\symidx{\(\operatorname{Cl}(R)\)：正规整环的 Weil 除子类群}

主除子的和确实有限。对 \(0\ne a\in R\)，只须证明：在高度一素理想中，\(v_{\mathfrak p}(a)>0\) 的那些恰好是 \((a)\) 的极小素理想。首先，DVR 中正阶恰表示属于极大理想，故
\(v_{\mathfrak p}(a)>0\) 当且仅当 \(a\in\mathfrak p\)。若 \(\mathfrak q\) 是 \((a)\) 的极小素理想，\cref{b3:thm:principal-ideal-theorem}给出 \(\operatorname{ht}\mathfrak q\le1\)；又因 \(R\) 是整环且 \(a\ne0\)，有 \((0)\subsetneq\mathfrak q\)，故其高度恰为一。反之，若高度一素理想 \(\mathfrak p\) 包含 \((a)\)，可取 \((a)\) 的极小素理想 \(\mathfrak q\subseteq\mathfrak p\)。两者都高度一，若包含严格，则 \((0)\subsetneq\mathfrak q\subsetneq\mathfrak p\) 与 \(\operatorname{ht}\mathfrak p=1\) 矛盾，故 \(\mathfrak p=\mathfrak q\)。Noether 性保证 \((a)\) 只有有限个极小素理想。一般非零分式 \(f=a/b\) 的非零阶数只可能出现在 \((a)\) 或 \((b)\) 的极小素理想处，因此支撑有限。赋值的可加性给出
\(\operatorname{div}(fg)=\operatorname{div}(f)+\operatorname{div}(g)\)，所以主除子确实构成子群。定义及其与 Picard 群的比较可参阅\cite[\S20, pp. 165--167]{Matsumura1989CommutativeRingTheory}。

\ProseParagraphBreak
一个基除子 \([\mathfrak p]\) 在类群中为零，意味着存在 \(f\in K^*\)，使它在 \(\mathfrak p\) 处恰有一阶零，在其他高度一处的阶数都为零。交定理将使 \(f\) 属于 \(R\)，并使 \(\mathfrak p=(f)\)。因此，这个类衡量高度一素理想能否由一个全局方程生成；下面的证明把这种解释与唯一分解判据联系起来。

\begin{theorem}\label{b3:thm:class-zero-ufd}
设 \(R\) 为交换 Noether 正规整环。则 \(R\) 为唯一分解整环，当且仅当 \(\operatorname{Cl}(R)=0\)。
\end{theorem}
\begin{proof}
若 \(R\) 是 UFD，则每个高度一素理想都是 \((\pi)\)，其中 \(\pi\) 为素元，并有 \(\operatorname{div}(\pi)=[(\pi)]\)。因此除子群的每个基元素均为主除子。

反之，若类群为零，给定高度一素理想 \(\mathfrak p\)，存在 \(f\in K^*\) 满足 \(\operatorname{div}(f)=[\mathfrak p]\)。所有赋值非负，所以 \(f\in R\)，且 \(f\in\mathfrak p\)。对任意非零的 \(a\in\mathfrak p\)，在 \(\mathfrak p\) 处有 \(v_{\mathfrak p}(a/f)\ge0\)，在其他高度一素理想处同样非负。交定理给出 \(a/f\in R\)，而零元也属于 \((f)\)，故 \(\mathfrak p=(f)\)。

这样每个高度一素理想都为主理想，\cref{b3:prop:height-one-ufd-criterion}对当前 Noether 整环适用，给出 \(R\) 为 UFD。
\end{proof}

\subsection{Dedekind 理论的高维延伸}
\label{b3:subsec:2003:03}

在 Dedekind 整环中，每个非零素理想都高度一，非零分式理想按赋值唯一分解。这时除子 \(\sum n_{\mathfrak p}[\mathfrak p]\) 对应分式理想 \(\prod\mathfrak p^{n_{\mathfrak p}}\)，本章的类群与\secref{b3:sec:1004}的理想类群一致。

\ProseParagraphBreak
在高维正规整环中，函数的上述恢复性质不能推广到任意理想。例如，设 \(k\) 为域，\(R=k[x,y]\)。理想 \(I=(x,y)\) 在每个高度一素点处均变为整个局部环：高度一素理想不可能同时包含 \(x,y\)。所以 \(I\) 与 \(R\) 有相同的高度一局部数据，但 \(I\ne R\)。特别地，在分式域中有
\[
\bigcap_{\operatorname{ht}\mathfrak p=1}I_{\mathfrak p}
=\bigcap_{\operatorname{ht}\mathfrak p=1}R_{\mathfrak p}
=R\supsetneq I.
\]
函数可以由高度一处的无极点条件恢复，这个普通理想在原点要求的同时消失却无法由那些条件恢复。除子记录的是余维一信息，不记录这个集中在余维二闭点的差别。

\ProseParagraphBreak
要把 Dedekind 分式理想的群运算推广到高维，就应选取恰能由高度一数据恢复的一类模。\secref{b3:sec:2005}将证明，秩一反射模具有这一性质；普通理想乘积则要再取双对偶，才能回到这类模中。

\ProseParagraphBreak
局部化则直接改变哪些高度一素理想仍需记录。逆转一个乘法集后，与它相交的素理想不再对应局部化环的素点；除子群删去相应基元素，类群便按这些类取商。

\begin{proposition}\label{b3:prop:class-group-localization}
设 \(R\) 为交换 Noether 正规整环，\(K=\operatorname{Frac}(R)\)，\(S\subseteq R\setminus\{0\}\) 为乘法集，\(A=S^{-1}R\)。则 \(A\) 仍为 Noether 正规整环，且保留避开 \(S\) 的高度一素理想所诱导的自然群同态
\[
\operatorname{Cl}(R)\longrightarrow\operatorname{Cl}(A)
\]
满射，其核为满足 \(\mathfrak p\cap S\ne\varnothing\) 的高度一素理想 \(\mathfrak p\) 的类所生成的子群。
\end{proposition}
\begin{proof}
局部化保持 Noether 性及整性。\(A\) 的每个素点局部环都同构于 \(R\) 的相应素点局部环，所以正规性按定义保持；两环的分式域都为 \(K\)。素理想对应把 \(\operatorname{Spec}A\) 与 \(R\) 中避开 \(S\) 的素理想对应起来。若 \(\mathfrak p\cap S=\varnothing\)，则 \(\mathfrak p\) 以下的每个素理想也避开 \(S\)，故对应保持止于 \(\mathfrak p\) 的所有素链，给出
\[
\operatorname{ht}_A(S^{-1}\mathfrak p)=\operatorname{ht}_R\mathfrak p.
\]
因此，两边的高度一素理想恰按这个对应匹配，并且
\(A_{S^{-1}\mathfrak p}\cong R_{\mathfrak p}\)，作为 \(K\) 的子环相同。相应归一化赋值也就相同。

在除子群上定义 \(\lambda\)：将避开 \(S\) 的 \([\mathfrak p]\) 送到 \([S^{-1}\mathfrak p]\)，将与 \(S\) 相交的 \([\mathfrak p]\) 送到零。每个 \(A\) 的除子基元素都有这样的原像，所以 \(\lambda\) 满射；其核恰由被删去的基元素生成。赋值相同又使
\[
\lambda\bigl(\operatorname{div}_R(f)\bigr)
=\operatorname{div}_A(f)\qquad(f\in K^*).
\]
这说明 \(\lambda\) 诱导所述类群满射。

若除子 \(D\) 的类落在这个满射的核中，则存在 \(f\in K^*\)，使 \(\lambda(D)=\operatorname{div}_A(f)\)。于是
\(\lambda(D-\operatorname{div}_R(f))=0\)，这个差除子只支撑在与 \(S\) 相交的高度一素理想上。它与 \(D\) 在 \(\operatorname{Cl}(R)\) 中给出同一个类，因此核包含于所述生成子群。反向，这些基除子都被 \(\lambda\) 删除，它们的类均在核中。
\end{proof}
